\chapter{齿轮机构 (Gear Mechanisms)}

\intro{齿轮机构是现代机械中应用最广泛的传动机构，用于传递任意两轴之间的运动和动力。其核心优势在于传动比准确、效率高、寿命长。}

\section{齿轮机构的应用及其分类}

齿轮机构根据两轴相对位置的不同，可分为以下几类：

\begin{enumerate}
	\item \textbf{平行轴传动}：直齿圆柱齿轮、斜齿圆柱齿轮、人字齿轮。
	\item \textbf{相交轴传动}：锥齿轮（直齿、斜齿、曲齿）。
	\item \textbf{交错轴传动}：蜗杆传动、交错轴斜齿轮。
\end{enumerate}

\section{齿廓啮合基本定律}

\begin{definition}[传动比(drive ratio)]
	两轮的瞬时角速度之比。
\end{definition}
\begin{definition}[齿廓啮合基本定律]
	互相啮合的一对齿轮，在任一时刻的传动比，等于两轮连心线被齿廓接触点的公法线所分成的两线段的反比。若传动比为常数，则节点为定点。
\end{definition}

\begin{definition}[节点(pitch point)]
	两齿轮瞬心。
\end{definition}

\begin{definition}[瞬心线(instantaneous center line)]
	节点$P$在每个齿轮运动平面上的轨迹。
\end{definition}

\begin{definition}[节圆(pitch circle)]
	圆形齿轮节点在两个齿轮运动平面上的轨迹是两个圆。
\end{definition}


\begin{definition}[共轭齿廓(conjugate profiles)]
	满足齿廓啮合基本定律的一对齿轮的齿廓。
\end{definition}
\subsection{数学推导}
设两齿轮的角速度分别为 $\omega_1, \omega_2$，接触点为 $K$，过 $K$ 点作两齿廓的公法线 $nn$，与连心线 $O_1O_2$ 交于 $P$ 点（节点）。
根据运动学速度合成定理，接触点 $K$ 分别在两构件上的速度 $v_{K1}$ 和 $v_{K2}$ 在公法线方向的分量必须相等：
\begin{equation}
	v_{K1} \cos \alpha = v_{K2} \cos \beta
\end{equation}
通过几何几何关系推导，可得传动比方程：
\begin{equation}
	i_{12} = \frac{\omega_1}{\omega_2} = \frac{O_2P}{O_1P}
\end{equation}
若要使传动比为常数，则节点 $P$ 在连心线上的位置必须固定。

\section{渐开线形成及其特性}

\subsection{渐开线的形成}
\begin{definition}{渐开线 (Involute)}
	一条直线（发生线）沿一个圆（基圆）作纯滚动时，直线上任一点的轨迹。
\end{definition}

\subsection{渐开线的特性}
\begin{enumerate}
	\item 发生线在基圆上的滚动长度等于基圆上对应的弧长：$BK = \wideparen{AB}$。
	\item 渐开线上任意一点的公法线必切于基圆。且越接近基圆的点曲率半径越小，曲率越大、
	\item 离基圆越远，压力角越大。基圆内无渐开线。
\end{enumerate}

\subsection{渐开线函数 (Involute Function)}
设渐开线上任意一点的向径为 $r$，压力角为 $\alpha$，基圆半径为 $r_b$。由渐开线几何基本关系：
\[
r \cos\alpha = r_b
\]

根据渐开线形成原理，渐开线展角 $\theta$ 与压力角 $\alpha$ 满足几何关系：
\[
\theta = \tan\alpha - \alpha
\]

定义**渐开线函数**为展角 $\theta$，记作 $\operatorname{inv}\alpha$，于是得到渐开线函数标准方程：
\begin{equation}
	\operatorname{inv} \alpha = \tan \alpha - \alpha
\end{equation}
式中角度 $\alpha$ 须以**弧度**为单位。渐开线函数为超越函数，是齿轮齿形几何参数计算、变位系数求解中最核心的基础方程。

以基圆半径 $r_b$ 和展角 $\theta$ 为参数，渐开线的参数方程可表示为：
\[
\begin{cases}
	x = r_b(\cos\theta + \theta\sin\theta)\\
	y = r_b(\sin\theta - \theta\cos\theta)
\end{cases}
\]

利用平面曲线曲率公式，可求得渐开线上任意点的**曲率半径**：
\[
\rho = r_b \theta
\]

结合渐开线函数关系，曲率半径也可通过压力角直接计算：
\[
\rho = r_b (\tan\alpha - \alpha)
\]
该式表明，渐开线的曲率半径随压力角增大而线性增加，是齿轮啮合分析与齿面强度计算的重要几何参数。

\section{渐开线齿廓啮合的特点}
\begin{enumerate}
	\item 啮合线是直线
	\item 啮合角不变
	\item 具有可分性
\end{enumerate}
\begin{proposition}
	如何理解渐开线齿廓啮合的可分性？
\end{proposition}
渐开线齿轮的传动比取决于其基圆的大小（简单画个图，利用相似），而一对齿轮加工好之后基圆大小不发生改变。


\section{渐开线齿轮各部分名称及标准齿轮的尺寸}

\subsection{渐开线齿轮各部分名称（Nomenclature of Involute Gear）}
以标准渐开线直齿圆柱齿轮为例，其轮齿各几何组成部分中英文名称及定义如下：

\begin{enumerate}
	\item \textbf{齿顶圆（Addendum Circle）}：齿轮齿顶所在的圆周，半径记为 $r_a$，直径记为 $d_a$；
	\item \textbf{齿根圆（Dedendum Circle）}：齿轮齿槽底部所在的圆周，半径记为 $r_f$，直径记为 $d_f$；
	\item \textbf{分度圆（Reference Circle / Pitch Circle）}：标准齿轮设计基准圆，半径 $r$、直径 $d$，为齿轮几何计算与啮合传动的基准；
	\item \textbf{基圆（Base Circle）}：生成渐开线齿廓的基准圆，半径 $r_b$、直径 $d_b$，渐开线由基圆渐开展成；
	\item \textbf{齿顶高（Addendum）}：分度圆至齿顶圆的径向高度，记作 $h_a$；
	\item \textbf{齿根高（Dedendum）}：分度圆至齿根圆的径向高度，记作 $h_f$；
	\item \textbf{全齿高（Whole Tooth Height）}：齿顶圆到齿根圆的径向总高度，满足 $h = h_a + h_f$；
	\item \textbf{齿厚（Tooth Thickness）}：分度圆上同一轮齿两侧齿廓间的弧长，记作 $s$；
	\item \textbf{齿槽宽（Space Width）}：分度圆上相邻两齿间齿槽的弧长，记作 $e$；
	\item \textbf{齿距（Circular Pitch）}：分度圆上相邻两齿对应点的弧长，满足 $p = s + e$。
\end{enumerate}
\subsection{基本参数}
\begin{itemize}
	\item \textbf{齿数 $z$}：齿轮上的轮齿总数。
	\item \textbf{模数 $m$}：$m = \frac{p}{\pi}= \frac{d}{z}$，其中 $p$ 为齿距。国家标准系列值。
	\item \textbf{分度圆压力角 $\alpha$}：标准值为 $20^\circ$。
\end{itemize}

\subsection{标准齿轮几何尺寸计算表}
\begin{table}[H]
	\centering
	\color{structurecolor}
	\begin{tabular}{lll}
		\hline
		\textbf{名称} & \textbf{代号} & \textbf{计算公式} \\ \hline
		分度圆直径 & $d$ & $d = mz$ \\
		分度圆齿距 & $p$ & $p = \pi m$ \\
		分度圆齿厚 & $s$ & $s = \frac{\pi m}{2} $ \\
		分度圆齿槽宽 & $e$ & $e = \frac{\pi m}{2} $ \\
		基圆直径 & $r_b$ & $d_b = d \cos \alpha$ \\
		齿顶高系数 & $h_a^*$ & 一般取 $1.0$  短齿取$0.8$\\
		齿顶高 &$h_a$ & $h_f=h_a^* m$ \\
		顶隙系数 & $c^*$ & 一般取 $0.25$  短齿取$0.3$ \\
		齿根高 &$h_f$ & $h_f=(h_a^*+c^*) m$ \\
		齿全高 &$h$ & $h=(2h_a^*+c^*) m$ \\
		齿顶圆直径 & $d_a$ & $d_a = d + 2h_a^* m$ \\
		齿根圆直径 & $d_f$ & $d_f = d - 2(h_a^* + c^*)m$ \\ \hline
	\end{tabular}
\end{table}

\section{渐开线直齿圆柱齿轮的啮合传动}

\subsection{正确啮合条件}
两齿轮要进入正常啮合状态，其模数和压力角必须分别相等：
\begin{equation}
	m_1 = m_2 = m, \quad \alpha_1 = \alpha_2 = \alpha
\end{equation}

\subsection{标准齿轮传动的正确安装中心距}
\begin{equation}
	a = r_1 +r_2
\end{equation}
非正确安装时：
\begin{equation}
	a' = r'_1 +r'_2
\end{equation}

存在一个这样的关系“
\begin{equation}
	a'=a\frac{\cos \alpha}{\cos \alpha'}
\end{equation}
\subsection{重合度 (Contact Ratio)}
为保证连续传动，实际啮合线段 $B_1B_2$ 必须大于法向齿距 $p_b$：
\begin{equation}
	\epsilon_\alpha = \frac{1}{2\pi}[z_1(\tan \alpha_{a1}-\tan \alpha')+z_2(\tan \alpha_{a2}-\tan \alpha')]
\end{equation}

\section{渐开线齿轮的加工与根切}

\subsection{加工方法}
主要分为成形法（铣齿）和展成法（滚齿、插齿）。

\subsection{根切现象 (Undercutting)}
当用展成法加工齿数过少的齿轮时，刀具齿顶会切去轮齿根部的渐开线齿廓。
\begin{theorem}{不发生根切的最小齿数}
	对于标准齿轮（$\alpha = 20^\circ, h_a^* = 1$），其不发生根切的最小齿数为：
	\begin{equation}
		z_{min} = \frac{2h_a^*}{\sin^2 \alpha} \approx 17
	\end{equation}
\end{theorem}

\section{变位齿轮 (Modified Gears)}

为了避免根切或凑中心距，需要将刀具相对于齿轮中心径向移动一段距离 $xm$。
\begin{itemize}
	\item $x$：变位系数。
	\item \textbf{正变位} ($x>0$)：刀具远离中心，齿根变厚，强度提高。
	\item \textbf{负变位} ($x<0$)：刀具靠近中心，易根切。
\end{itemize}

\section{斜齿圆柱齿轮机构}

\subsection{斜齿轮的特点}
1. 啮合性能好，传动平稳。
2. 重合度大。
3. 产生轴向力。

\subsection{当量齿数 (Equivalent Number of Teeth)}
由于斜齿轮法面齿廓为渐开线，通过法面代换，其当量齿数为：
\begin{equation}
	z_v = \frac{z}{\cos^3 \beta}
\end{equation}
其中 $\beta$ 为螺旋角。

\section{蜗杆传动 (Worm Drive)}

\subsection{传动特点}
\begin{enumerate}
	\item 传动比大（通常 $i = 10 \sim 80$）。
	\item 结构紧凑。
	\item 具有反向自锁性。
	\item 效率低，发热量大（存在剧烈滑动）。
\end{enumerate}

\begin{solution}
	\textbf{例题}：已知一标准直齿圆柱齿轮，$z=20, m=4mm, \alpha=20^\circ$，求其基圆半径 $r_b$。
	
	解：首先计算分度圆半径：
	$r = \frac{mz}{2} = \frac{4 \times 20}{2} = 40 mm$
	
	根据基圆与分度圆关系：
	$r_b = r \cos \alpha = 40 \times \cos 20^\circ \approx 37.588 mm$
\end{solution}